% part: intuitionistic-logic
% chapter: introduction
% section: bhk-interpretation

\documentclass[../../../include/open-logic-section]{subfiles}

\begin{document}

\olfileid{int}{int}{bhk}

\olsection{د براور--هايتنګ--کولموګوروف تفسير}

\begin{editorial}
  د BHK تفسير په کارولو سره د شهودي قضيو د اعتبار ثبوتونه مغلق دي؛
  بايد لا ښه تشريح شي۔
\end{editorial}

د شهودي نښلوونکو يو غيررسمي تعميري تفسير شته چې عموماً د
براور--هايتنګ--کولموګوروف تفسير بلل کېږي۔ دا د «جوړونې» مفهوم کاروي
چې د تعميري ثبوت په توګه ئې تصورولے شئ۔ (په BHK تفسير کښې د «ثبوت»
پر ځاے «جوړونه» وايو، څو په رسمي !!{derivation} نظام کښې له
!!a{derivation} سره ګډ نۀ شي۔) د همدې شهودي مفهوم په بنسټ BHK تفسير
د شهودي نښلوونکو معناوې بيانوي۔

\begin{enumerate}
\item فرض کوو چې د يوې اټومي جملې جوړونه څۀ شی دے، دا موږ ته معلومه ده۔
\item د $!A_1 \land !A_2$ جوړونه يوه جوړه $\tuple{M_1, M_2}$ ده
  چې $M_1$ د~$!A_1$ جوړونه او $M_2$ د~$!A_2$ جوړونه وي۔
\item د $!A_1 \lor !A_2$ جوړونه يوه جوړه $\tuple{s, M}$ ده
  چې يا $s$ له~$1$ سره برابر وي او $M$ د~$!A_1$ جوړونه وي، يا $s$
  له~$2$ سره برابر وي او $M$ د~$!A_2$ جوړونه وي۔
\item د $!A \lif !B$ جوړونه يو تابع دے چې د~$!A$ جوړونه
  د~$!B$ په جوړونه اړوي۔
\item د~$\lfalse$ (محال) هېڅ جوړونه نۀ شته۔
\item $\lnot !A$ د $!A \lif \lfalse$ مترادفه نښه تعريفېږي۔ يعنې
  د $\lnot !A$ جوړونه يو تابع دے چې د~$!A$ جوړونه د~$\lfalse$
  په جوړونه اړوي۔
\end{enumerate}

\begin{ex}
د بېلګې په توګه $\lnot \lfalse$ واخلئ۔ د دې جوړونه داسې تابع دے
چې د $\lfalse$ هره ورکړل شوې جوړونه د~$\lfalse$ په جوړونه اړوي۔
ښکاره ده چې عيني تابع~$\Id{}$ داسې جوړونه ده: که د~$\lfalse$
جوړونه~$M$ ورکړل شي، $\Id{}(M) = M$ بيا د~$\lfalse$ جوړونه ورکوي۔
\end{ex}

په عمومي ډول، $\lnot !A$ دا معنا لري چې «د $!A$ جوړونه ناشونې ده»۔

\begin{ex}
د هرې قضيې~$!A$ لپاره $!A \lif \lnot \lnot !A$ ثابتوو، چې
$!A \lif ((!A \lif \lfalse) \lif \lfalse)$ دے۔ جوړونه بايد داسې
تابع~$f$ وي چې د $!A$ يوه جوړونه~$M$ واخلي او د
$(!A \lif \lfalse) \lif \lfalse$ يوه جوړونه~$f(M)$ ورکړي۔
تابع~$f$ د $(!A \lif \lfalse) \lif \lfalse$ جوړونه داسې جوړوي:
بايد داسې تابع $g$ تعريف کړو چې د $!A \lif \lfalse$ جوړونه~$h$
واخلي او د~$\lfalse$ جوړونه ورکړي۔ $g$ داسې تعريفولے شو: ورکړل شوے
$h$ د $!A$ پر هغې جوړونې~$M$ تطبيق کړئ چې مخکې مو اخيستې وه۔ ځکه د
$h$ پايله~$h(M)$ د $\lfalse$ جوړونه ده، نو که $M$ د~$!A$ جوړونه
وي، $f(M)(h) = h(M)$ د~$\lfalse$ جوړونه ده۔
\end{ex}

\begin{ex}
د $\lnot(!A \land \lnot !A)$ لپاره يوه جوړونه ورکوو، يعنې د $(!A
\land (!A \lif \lfalse)) \lif \lfalse$ لپاره۔ دا داسې تابع~$f$
دے چې د $!A \land (!A \lif \lfalse)$ جوړونه~$M$ واخلي او د
~$\lfalse$ جوړونه ورکړي۔ د عطف $!B_1 \land !B_2$ جوړونه يوه
جوړه $\tuple{N_1, N_2}$ ده چې $N_1$ د~$!B_1$ او $N_2$ د~$!B_2$
جوړونه وي۔ موږ $p_1$ او $p_2$ داسې تابعې تعريفولے شو چې د
$!B_1 \land !B_2$ له جوړونې څخه په ترتيب د $!B_1$ او $!B_2$
جوړونې راوباسي:
\begin{align*}
  p_1(\tuple{N_1, N_2}) &= N_1\\
  p_2(\tuple{N_1, N_2}) &= N_2
\end{align*}
تابع $f$ داسې کار کوي: لومړے $p_1$ په خپل ورودي~$M$ تطبيقوي۔
په دې سره د~$!A$ جوړونه لاس ته راځي۔ بيا $p_2$ په~$M$
تطبيقوي او د~$!A \lif \lfalse$ جوړونه ترې لاس ته راځي۔ دا
جوړونه پخپله يو تابع~$p_2(M)$ دے چې د~$!A$ جوړونه واخلي او د
~$\lfalse$ جوړونه ورکړي۔ په بله وينا، که $p_2(M)$ په $p_1(M)$
تطبيق کړو، د~$\lfalse$ جوړونه تر لاسه کوو۔ نو
$f(M) = p_2(M)(p_1(M))$ تعريفولے شو۔
\end{ex}

\begin{ex}
د $((!A \land !B) \lif !C) \lif (!A \lif
(!B \lif !C))$ جوړونه ورکوو۔ يعنې داسې تابع $f$ ورکوو چې د
$(!A \land !B) \lif !C$ جوړونه~$g$ د $(!A \lif (!B \lif
!C))$ په جوړونه واړوي۔ جوړونه $g$ پخپله يو تابع دے: د
$!A \land !B$ جوړونې د~$!C$ په جوړونو اړوي۔ د $f(g)$ پايله
يو تابع~$h_g$ دے چې د $!A$ جوړونې د هغو تابعو په جوړونو اړوي
چې د~$!B$ جوړونې د~$!C$ په جوړونو اړوي۔
\textbf{د ژباړې د نسخې څرګندونه (OLINT-001):} په منجمده سرچينه کښې
دلته د C له مخې د فارمول نښه پرېوتې ده؛ د همدې جملې او قضیې په نورو
فورمولونو کښې شته، نو پښتو متن يې د پايلې سره ساتي۔

ښه، دا مغلق ښکاري۔ داسې تابع $h_g$ جوړول غواړو چې د ورودي~$g$
لپاره د $f$ پايله وي۔ د~$h_g$ ورودي د~$!A$ جوړونه~$M$ ده۔ د
$h_g(M)$ پايله بايد داسې تابع~$k_M$ وي چې د~$!B$ جوړونې~$N$
د~$!C$ په جوړونو اړوي۔ $k_{g,M}(N) = g(\tuple{M, N})$ وټاکئ۔
په ياد ولرئ چې $\tuple{M, N}$ د~$!A \land !B$ جوړونه ده۔ نو
$k_{g,M}$ د $!B \lif !C$ جوړونه ده: د~$!B$ جوړونې~$N$
د~$!C$ په جوړونو اړوي۔ اوس $h_g(M) = k_{g,M}$ وټاکئ۔ دا داسې
تابع دے چې د~$!A$ جوړونې~$M$ د $!B
\lif !C$ په جوړونو~$k_{g,M}$ اړوي۔ اوس $f(g) = h_g$ وټاکئ۔
دا داسې تابع ده چې د $(!A \land !B) \lif !C$ جوړونې~$g$
د $!A \lif (!B \lif !C)$ په جوړونو اړوي۔ اوس خبره څرګنده شوه!{}
\end{ex}

جمله $!A \lor \lnot !A$ د منځني حالت د ايستلو اصل بلل کېږي۔ د ځينو
ځانګړو $!A$ لپاره ئې ثابتولے شو، لکه $\lfalse \lor
\lnot \lfalse$، خو په عمومي ډول نه۔ وجه دا ده چې شهودي فصل د
دوو اړخونو څخه د يوه جوړونه غواړي، حال دا چې ځينې جملې اوس نۀ
ثابتېدے شي او نۀ رد کېدے شي، لکه د ګولډباخ حدسيه۔ خو د
منځني حالت د ايستلو اصل هم نۀ شئ ردولے: يعنې $\lnot \lnot (!A
\lor \lnot !A)$ رښتيا دے۔

\begin{ex}
د $\lnot \lnot (!A \lor \lnot !A)$ د ثبوت لپاره داسې تابع~$f$
غواړو چې د $\lnot(!A \lor \lnot !A)$ جوړونه، يعنې د $(!A
\lor (!A \lif \lfalse)) \lif \lfalse$ جوړونه، د~$\lfalse$
په جوړونه واړوي۔ په بله وينا، داسې تابع $f$ غواړو چې که $g$ د
$\lnot(!A \lor \lnot !A)$ جوړونه وي، نو $f(g)$ د~$\lfalse$
جوړونه وي۔

فرض کړئ $g$ د $\lnot(!A \lor \lnot !A)$ جوړونه ده؛ يعنې داسې
تابع چې د $!A \lor \lnot !A$ جوړونه د~$\lfalse$ په جوړونه
اړوي۔ د $!A \lor \lnot !A$ جوړونه يوه جوړه $\tuple{s, M}$ ده
چې يا $s = 1$ وي او $M$ د~$!A$ جوړونه وي، يا $s = 2$ وي او
$M$ د $\lnot !A$ جوړونه وي۔ $h_1$ هغه تابع وټاکئ چې د $!A$
جوړونه~$M_1$ د $!A \lor \lnot !A$ په جوړونه اړوي: $M_1$
د $\tuple{1, M_1}$ جوړې ته لېږي۔ $h_2$ هغه تابع وټاکئ چې د
$\lnot !A$ جوړونه~$M_2$ د $!A \lor \lnot !A$ په جوړونه
اړوي: $M_2$ د $\tuple{2, M_2}$ جوړې ته لېږي۔
\textbf{د ژباړې د نسخې څرګندونه (OLINT-002):} منجمده سرچينه د لومړي
تابع په جمله کښې ورکړل شوې لومړۍ جوړونه د دويمې جوړونې په نښه لرونکې
جوړې ته لېږي؛ د فصل د لومړۍ مادې له مخې د جوړې دويم غړی بايد هماغه
ورکړل شوې لومړۍ جوړونه وي۔

$k$ د $\comp{h_1}{g}$ په توګه وټاکئ: دا داسې تابع دے چې د~$!A$
جوړونه واخلي او د~$\lfalse$ جوړونه ورکړي؛ يعنې د
~$!A \lif \lfalse$ يا $\lnot !A$ جوړونه ده۔ اوس $l$ د
$\comp{h_2}{g}$ په توګه وټاکئ۔ دا داسې تابع ده چې د $\lnot !A$
جوړونه واخلي او د~$\lfalse$ جوړونه ورکړي۔ څرنګه چې $k$ د
$\lnot !A$ جوړونه ده، نو $l(k)$ د~$\lfalse$ جوړونه ده۔

په ګډه مو وښودله چې د $\lnot(!A \lor \lnot !A)$ جوړونه~$g$
څنګه د~$\lfalse$ په جوړونه اړولے شو۔ يعنې هغه تابع $f$ چې د
$\lnot(!A \lor \lnot !A)$ جوړونه~$g$ د~$\lfalse$ په
جوړونه $l(k)$ اړوي، د $\lnot\lnot(!A \lor \lnot !A)$ جوړونه ده۔
\end{ex}

لکه چې وينئ، د BHK تفسير په کارولو د !!{formula}s شهودي اعتبار
ښودل ژر اوږد او مغلق کېږي۔ له نېکه مرغه د شهودي منطق لپاره تر دې
غوره !!{derivation} نظامونه او لا دقيق معنايي تفسيرونه شته۔
\end{document}
